%%%%%%%%%%%%%%%%%%%%%%%%%%%%%%%%%%%%%%%%%%%%%%%%%%%%%%%%%%%%%%%%%%%
%%%
%%% File ini ditujukan untuk menulis
%%% Paper POMITS menggunakan LaTeX untuk program sarjana
%%% di Institut Teknologi Sepuluh Nopember, Surabaya.
%%%
%%%%%%%%%%%%%%%%%%%%%%%%%%%%%%%%%%%%%%%%%%%%%%%%%%%%%%%%%%%%%%%%%%%%
\documentclass[journal]{IEEEtran}
\hyphenation{op-tical net-works semi-conduc-tor}
\usepackage{amssymb,amsmath,amsthm,graphicx}
\usepackage{caption,tabularx,multicol,multirow}
\usepackage{listings,color,booktabs}
\usepackage{float}
\usepackage{algpseudocode}
\usepackage{array}
\usepackage{scrfontsizes}
\usepackage{afterpage}
\usepackage{subfig}
\usepackage{mathtools}
\usepackage{physics}
\usepackage[ruled,vlined]{algorithm2e}
\captionsetup[figure]{labelsep=period}
\captionsetup[figure]{name={Gambar}}
\captionsetup[table]{labelsep=newline}
\captionsetup[table]{name={Tabel}}
\tolerance=1
\emergencystretch=\maxdimen
\hyphenpenalty=10000
\hbadness=10000
\setlength{\parindent}{0.36cm}

\newtheorem{defn}{Definisi}[section]
\newtheorem{teo}[defn]{Teorema}
\newtheorem{thm}{Teorema}
\newtheorem{lemma}{Lemma}
\newtheorem{lemmas}[defn]{Lemma}
\newtheorem{cor}[defn]{Akibat}
\newtheorem{asumsi}[defn]{Asumsi}
\theoremstyle{definition}
\newtheorem{con}{Contoh}[section]
\theoremstyle{plain}
\newtheorem{prop}{Proposisi}
\renewcommand{\proofname}{Bukti}
\renewcommand{\thedefn}{\arabic{section}.\arabic{defn}}
\renewcommand{\thecon}{\arabic{section}.\arabic{con}}

\newcommand{\katakunci}[1]{%
    \noindent\\
    \textit{\textbf{Kata Kunci---}}{\textbf{#1}}
}

\renewcommand{\listfigurename}{\hfill\normalsize\textbf{DAFTAR GAMBAR}\hfill}
\renewcommand{\listtablename}{\hfill\normalsize\textbf{DAFTAR TABEL}\hfill}
\renewcommand{\contentsname}{\hfill\normalsize\textbf{DAFTAR ISI}\hfill}
\renewcommand{\refname}{}
\renewcommand{\arraystretch}{1.2}

\newcolumntype{L}[1]{>{\raggedright\let\newline\\\arraybackslash\hspace{0pt}}m{#1}}
\newcolumntype{C}[1]{>{\centering\let\newline\\\arraybackslash\hspace{0pt}}m{#1}}
\newcolumntype{R}[1]{>{\raggedleft\let\newline\\\arraybackslash\hspace{0pt}}m{#1}}

\definecolor{codegreen}{rgb}{0,0.6,0}
\definecolor{codegray}{rgb}{0.5,0.5,0.5}
\definecolor{codepurple}{rgb}{0.58,0,0.82}
\definecolor{ITSblue}{rgb}{0.08,0.38,0.74}

\lstdefinestyle{mypseudo}{
commentstyle=\bfseries,
keywordstyle=\bfseries,
breakatwhitespace=false,
breaklines=true,
captionpos=b,
keepspaces=true,
showspaces=false,
showstringspaces=false,
showtabs=false,
literate={\ \ }{{\ }}2
}
\lstset{basicstyle=\sffamily\scriptsize,style=mypseudo}
\def\abstractname{Abstrak}

\newcommand{\Rmax}{\mathbb{R}_{\max}}
\newcommand{\LS}{L_S(n)}
\newcommand{\Sn}{S_n}
\newcommand{\eps}{\varepsilon}

\begin{document}

\title{\changefontsizes{24}\textit{Latin Square} Komutatif atas Aljabar Max-Plus\\ \vspace{12pt}}

\author{\changefontsizes{11}Teosofi Hidayah Agung\\
Departemen Matematika, Fakultas Sains dan Analitika Data\\
Institut Teknologi Sepuluh Nopember (ITS), Surabaya, Indonesia\\ \vspace{-1em}}

\markboth{\small JURNAL TEKNIK ITS Vol. X, No. Y, (2026) ISSN: 2337-3539 (2301-9271 Print)}{}

\maketitle

\begin{abstract}
\small \noindent Aljabar max-plus adalah struktur aljabar pada himpunan
$\Rmax=\mathbb{R}\cup\{\eps\}$ dengan operasi
$a\oplus b=\max\{a,b\}$ dan $a\otimes b=a+b$.
Penelitian ini mengkaji komutativitas dua \textit{Latin square}
terhadap perkalian max-plus, yaitu kondisi
$A\otimes B=B\otimes A$.
Setiap \textit{Latin square} direpresentasikan sebagai kombinasi max-plus
dari matriks permutasi, sehingga komutativitas dapat dianalisis melalui
komposisi permutasi penyusunnya.
Hasil utama penelitian ini meliputi rumus dekomposisi hasil perkalian
\textit{Latin square}, syarat perlu komutativitas melalui simbol maksimum,
syarat cukup melalui subgrup abelian dari $\Sn$, serta kriteria perlu dan
cukup melalui kesamaan himpunan superlevel.
Berdasarkan kriteria tersebut disusun algoritma untuk mengonstruksi pasangan
\textit{Latin square} komutatif dari suatu \textit{Latin square} yang
diberikan.
Penerapan algoritma pada ordo $3$, $4$, dan $5$ menunjukkan bahwa prosedur
yang diperoleh dapat digunakan baik pada kasus ketika permutasi penyusun
berada dalam subgrup abelian yang sama maupun pada kasus umum melalui
\textit{centralizer} dan pemeriksaan superlevel.
\end{abstract}

\katakunci{\small Aljabar max-plus, \textit{Latin square}, komutativitas, \textit{centralizer}, himpunan superlevel.}

\IEEEpeerreviewmaketitle

\section{PENDAHULUAN}
\normalsize \IEEEPARstart{A}{LJABAR} max-plus merupakan salah satu struktur
dasar dalam aljabar tropikal dengan operasi penjumlahan berupa maksimum dan
operasi perkalian berupa penjumlahan biasa \cite{subiono2015minmaxplus}.
Struktur ini banyak digunakan dalam pemodelan sistem kejadian diskret,
optimasi, dan kajian kriptografi tropikal \cite{huang2024newsemiring,
durcheva2025cryptographyidempotent}.

\textit{Latin square} adalah matriks berordo $n$ yang setiap baris dan
kolomnya merupakan permutasi dari himpunan simbol yang sama
\cite{george2017introductioncombinatorics}.
Dalam konteks aljabar max-plus, \textit{Latin square} telah dikaji dari sisi
nilai eigen, vektor eigen, dan algoritma komputasi
\cite{mufid2014eigenvalues}.
Di sisi lain, matriks komutatif pada aljabar tropikal juga menarik karena
berkaitan dengan konstruksi protokol pertukaran kunci berbasis matriks
\cite{alhussaini2024securityinitialtropical,
alhussaini2024tropicaltwosided,muanalifah2020modifyingtropical,
buchinskiy2024analysisfour}. Kajian mengenai matriks yang komutatif dengan
matriks tropikal tertentu juga telah muncul, misalnya pada matriks normal
tropikal \cite{linde2014matricescommutinggivennormal}.

Berangkat dari latar belakang tersebut, penelitian ini berfokus pada tiga
persoalan, yaitu menentukan syarat perlu komutativitas dua
\textit{Latin square}, menentukan syarat cukup serta kriteria perlu-cukupnya,
dan menyusun algoritma untuk mengonstruksi pasangan \textit{Latin square}
komutatif.
Pembahasan dibatasi pada \textit{Latin square} dengan simbol
$\{1,2,\ldots,n\}$ dan penerapan algoritma diberikan untuk ordo $n\leq 5$.

\section{METODE PENELITIAN}
Metode penelitian diawali dengan studi literatur mengenai aljabar max-plus,
grup permutasi, matriks permutasi, \textit{centralizer}, dan
\textit{Latin square}. Selanjutnya dilakukan representasi
\textit{Latin square} ke dalam matriks-matriks permutasi max-plus, penurunan
kriteria komutativitas, penyusunan algoritma konstruksi pasangan komutatif,
dan pengujian hasil pada beberapa contoh berordo kecil.

Untuk $A=(a_{ij})$ dan $B=(b_{ij})$ di $\Rmax^{n\times n}$, perkalian
max-plus didefinisikan oleh
\begin{equation}
    [A\otimes B]_{rs}=\max_{1\leq t\leq n}(a_{rt}+b_{ts}).
\end{equation}

Jika $A\in \LS$ dengan himpunan simbol $\{1,2,\ldots,n\}$, maka posisi
kemunculan simbol $i$ pada setiap baris membentuk suatu permutasi
$\sigma_i^A\in \Sn$. Dengan demikian, setiap \textit{Latin square}
dapat ditulis sebagai
\begin{equation}
    A=\bigoplus_{i=1}^{n} i\otimes P_{\sigma_i^A}.
\end{equation}
Representasi yang sama berlaku untuk $B\in \LS$.

Dari representasi tersebut diperoleh dekomposisi hasil perkalian berikut.
\begin{cor}
Misalkan $A,B\in \LS$ dengan dekomposisi
$A=\bigoplus_{i=1}^{n} i\otimes P_{\sigma_i^A}$ dan
$B=\bigoplus_{j=1}^{n} j\otimes P_{\sigma_j^B}$.
Maka
\begin{equation}
    A\otimes B=
    \bigoplus_{k=n+1}^{2n}
    k\otimes
    \left(
    \bigoplus_{\substack{i,j\in\{1,\ldots,n\}\\i+j=k}}
    P_{\sigma_j^B\circ\sigma_i^A}
    \right).
\end{equation}
\end{cor}

Rumus tersebut menunjukkan bahwa komutativitas dua \textit{Latin square}
ditentukan oleh hubungan antara komposisi
$\sigma_j^B\circ\sigma_i^A$ dan $\sigma_j^A\circ\sigma_i^B$.
Atas dasar itu, penelitian dilanjutkan dengan menurunkan syarat perlu,
syarat cukup, dan kriteria perlu-cukup yang kemudian digunakan untuk
menyusun algoritma konstruksi pasangan komutatif.

\section{HASIL DAN PEMBAHASAN}
\subsection{Syarat Perlu Komutativitas}
Syarat perlu pertama diperoleh dari simbol maksimum pada
\textit{Latin square}.

\begin{teo}
Misalkan $A,B\in \LS$ dengan himpunan simbol $\{1,2,\ldots,n\}$.
Jika $A\otimes B=B\otimes A$, maka permutasi yang bersesuaian dengan simbol
maksimum memenuhi
\begin{equation}
    \sigma_n^A\circ\sigma_n^B=\sigma_n^B\circ\sigma_n^A.
\end{equation}
\end{teo}

Teorema ini memberikan penyaring awal pada proses pencarian kandidat.
Artinya, kandidat $\sigma_n^B$ hanya perlu dipilih dari
$C_{\Sn}(\sigma_n^A)$, yaitu \textit{centralizer} dari $\sigma_n^A$ di
$\Sn$.

\subsection{Syarat Cukup Komutativitas}
Selain syarat perlu, diperoleh pula syarat cukup berbasis subgrup abelian.

\begin{teo}
Misalkan $A,B\in \LS$ dengan dekomposisi
$A=\bigoplus_{i=1}^{n} i\otimes P_{\sigma_i^A}$ dan
$B=\bigoplus_{j=1}^{n} j\otimes P_{\sigma_j^B}$.
Jika terdapat subgrup abelian $H\leq \Sn$ sedemikian sehingga
$\sigma_i^A\in H$ dan $\sigma_j^B\in H$ untuk setiap
$i,j\in\{1,\ldots,n\}$, maka
\begin{equation}
    A\otimes B=B\otimes A.
\end{equation}
\end{teo}

Hasil ini penting karena langsung menghasilkan keluarga pasangan
\textit{Latin square} komutatif. Sebagai contoh, dua \textit{Latin square}
sirkulan selalu komutatif karena seluruh permutasi penyusunnya merupakan
pangkat-pangkat dari generator siklik yang sama.

\subsection{Kriteria Perlu dan Cukup}
Untuk memperoleh kriteria yang lengkap, digunakan himpunan superlevel.
Untuk setiap $v\in\{n+1,n+2,\ldots,2n\}$, didefinisikan
\begin{equation}
    \mathcal{U}_{\geq v}^{AB}=
    \bigcup_{\substack{i,j\in\{1,\ldots,n\}\\i+j\geq v}}
    \left\{(r,(\sigma_j^B\circ\sigma_i^A)(r))\mid r\in\{1,\ldots,n\}\right\}.
\end{equation}
Secara analog, didefinisikan pula $\mathcal{U}_{\geq v}^{BA}$.

\begin{teo}
Misalkan $A,B\in \LS$ dengan himpunan simbol $\{1,2,\ldots,n\}$.
Matriks $A$ dan $B$ komutatif terhadap perkalian max-plus jika dan hanya jika
\begin{equation}
    \mathcal{U}_{\geq v}^{AB}=\mathcal{U}_{\geq v}^{BA}
\end{equation}
untuk setiap $v\in\{n+2,n+3,\ldots,2n\}$.
\end{teo}

Kriteria ini mengubah persoalan kesamaan hasil kali matriks menjadi
pemeriksaan kesamaan himpunan posisi pada setiap level ambang.
Keuntungan utamanya adalah pemeriksaan dapat dilakukan secara bertahap dari
level tertinggi ke level yang lebih rendah.

\subsection{Algoritma Konstruksi Pasangan Komutatif}
Berdasarkan hasil sebelumnya, disusun algoritma untuk mencari semua kandidat
$B\in \LS$ yang memenuhi $A\otimes B=B\otimes A$ untuk suatu
$A\in \LS$ yang diberikan.
Secara ringkas, langkah-langkah algoritma adalah sebagai berikut.
\begin{enumerate}
    \item Dekomposisikan $A$ menjadi permutasi-permutasi penyusun
    $\sigma_1^A,\ldots,\sigma_n^A$.
    \item Jika seluruh permutasi penyusun $A$ saling komutatif, bentuk
    kandidat $B$ dengan menyusun ulang permutasi-permutasi tersebut.
    \item Pada tahap umum, pilih $\sigma_n^B$ dari
    $C_{\Sn}(\sigma_n^A)$ sebagai kandidat simbol maksimum.
    \item Bangun matriks parsial $B^*$ dan pilih $\sigma_m^B$ secara menurun
    untuk $m=n-1,n-2,\ldots,2$.
    \item Setiap kali sebuah permutasi dipilih, periksa apakah
    $\mathcal{U}_{\geq n+m}^{AB}=\mathcal{U}_{\geq n+m}^{BA}$.
    Jika tidak, cabang pencarian langsung ditolak.
    \item Setelah seluruh simbol selain $1$ terpasang, isi entri tersisa
    dengan simbol $1$ dan verifikasi bahwa posisinya membentuk matriks
    permutasi.
\end{enumerate}

Strategi ini memungkinkan pencarian dihentikan lebih awal ketika suatu
kandidat gagal pada level tertentu, sehingga tidak perlu melengkapi seluruh
permutasi penyusun $B$.

\subsection{Contoh Penerapan}
Berikut diberikan contoh dari tahap alternatif dan tahap umum pada algoritma.

\begin{con}
Misalkan
\begin{equation}
    A=
    \begin{pmatrix}
        1 & 2 & 3\\
        3 & 1 & 2\\
        2 & 3 & 1
    \end{pmatrix}.
\end{equation}
Permutasi penyusun $A$ saling komutatif, sehingga tahap alternatif dapat
digunakan. Salah satu pasangan komutatif yang diperoleh adalah
\begin{equation}
    B=
    \begin{pmatrix}
        1 & 3 & 2\\
        2 & 1 & 3\\
        3 & 2 & 1
    \end{pmatrix}.
\end{equation}
Perhitungan langsung memberikan
\begin{equation}
    A\otimes B=B\otimes A=
    \begin{pmatrix}
        6 & 5 & 5\\
        5 & 6 & 5\\
        5 & 5 & 6
    \end{pmatrix}.
\end{equation}
\end{con}

\begin{con}
Misalkan
\begin{equation}
    A=
    \begin{pmatrix}
        1 & 2 & 4 & 3\\
        3 & 1 & 2 & 4\\
        2 & 4 & 3 & 1\\
        4 & 3 & 1 & 2
    \end{pmatrix}.
\end{equation}
Permutasi penyusun $A$ tidak saling komutatif, sehingga pencarian dilakukan
melalui \textit{centralizer} dan pemeriksaan superlevel. Salah satu pasangan
yang diperoleh adalah
\begin{equation}
    B=
    \begin{pmatrix}
        4 & 1 & 3 & 2\\
        1 & 4 & 2 & 3\\
        2 & 3 & 4 & 1\\
        3 & 2 & 1 & 4
    \end{pmatrix}.
\end{equation}
Hasil perkaliannya adalah
\begin{equation}
    A\otimes B=B\otimes A=
    \begin{pmatrix}
        6 & 7 & 8 & 7\\
        7 & 6 & 6 & 8\\
        6 & 8 & 7 & 7\\
        8 & 7 & 7 & 6
    \end{pmatrix}.
\end{equation}
\end{con}

Penerapan pada ordo $5$ menunjukkan pola yang serupa. Pada contoh sirkulan,
tahap alternatif menghasilkan pasangan komutatif dari penyusunan ulang
permutasi penyusun. Pada contoh nonsirkulan, tahap umum melalui
\textit{centralizer} dan kesamaan superlevel juga menghasilkan kandidat yang
memenuhi $A\otimes B=B\otimes A$.

\section{KESIMPULAN}
Penelitian ini menunjukkan bahwa komutativitas dua \textit{Latin square}
pada aljabar max-plus dapat dipelajari secara sistematis melalui representasi
permutasi penyusunnya. Syarat perlu diperoleh dari simbol maksimum, syarat
cukup diperoleh dari keberadaan subgrup abelian yang memuat seluruh
permutasi penyusun, dan kriteria perlu-cukup diperoleh dari kesamaan
himpunan superlevel pada setiap level yang relevan. Berdasarkan hasil
tersebut, disusun algoritma konstruksi pasangan komutatif yang dapat
diterapkan baik pada kasus abelian maupun kasus umum.

Secara praktis, hasil ini memberikan prosedur yang lebih terarah untuk
membangun matriks komutatif berbentuk \textit{Latin square} dalam aljabar
max-plus. Penerapan pada ordo $3$, $4$, dan $5$ menunjukkan bahwa algoritma
yang diperoleh dapat digunakan untuk menghasilkan pasangan yang memenuhi
$A\otimes B=B\otimes A$.

\section*{UCAPAN TERIMA KASIH}
Penulis mengucapkan terima kasih kepada Muhammad Syifa'ul Mufid,
S.Si., M.Si., D.Phil. selaku dosen pembimbing yang telah memberikan arahan
dan masukan selama penyusunan penelitian ini.

\section*{DAFTAR PUSTAKA}
\bibliographystyle{IEEEtran}
\bibliography{file/daftarPustaka}

\end{document}
